RRF(3) уступает E5 при $K=10$ (0.852 против 0.871), но превосходит его при всех четырёх больших исследованных глубинах. При $K=200$ парная разность равна $+0.0088$, 95\%-й ДИ $[0.0082, 0.0095]$. Для всех пяти сравнений нескорректированное парное $p<0.001$; они также проходят поправку Бонферрони для пяти проверок при уровне 0.05. Это согласие с фиксированным судьёй, а не независимо подтверждённое превосходство по человеческим оценкам.

DCG первых десяти результатов отдельных поисковых моделей не меняется при росте $K$: изменяется лишь общий идеальный знаменатель. Более глубокие входные списки RRF, напротив, могут менять само итоговое ранжирование. При фиксированной глубине входа RRF, равной 200, NDCG составляет 0.805 при глубине разметки десять и 0.850 при глубине 200. В первом случае 145\,643 возвращённых пар находятся вне размеченного пула и получают нулевой вклад. Эта проверка отделяет неполноту разметки от усечения входов слияния.

Все 1\,179 запросов ручной части входят в большой набор, поэтому он не является полностью независимой тестовой выборкой. Их исключение оставляет 99\,914 запросов. При глубине 200 разность RRF минус E5 сохраняется: $+0.0088$ (95\%-й ДИ $[0.0080, 0.0094]$); направления всех пяти сравнений также не меняются. Проверены оценки для всех 45\,647\,587 пар без пропусков; при глубине 200 у 1\,592 запросов нет положительного предсказания.
